% title.tex
\title{Кинетика осуществления объектов весомой материи}
\author{Проект <<Эфир>>}
\date{}
\maketitle

\begin{abstract}
Наблюдаемая Вселенная состоит из соударяющихся частиц эфира, и основанием для такого утверждения является существование физических сценариев осуществления не разрушающихся со временем пространственных эфирных локаций — \textbf{объектов весомой материи}. К ним относятся все объекты наблюдаемой Вселенной, включая элементарные частицы и фотоны. \textbf{Весомость} — свойство, присущее организованно движущимся в ограниченном пространстве совокупностям частиц эфира, являющееся следствием физических процессов, реализующих непрерывную во времени пространственную локализацию объекта. В основе осуществления такой локализации лежат абсолютно упругие соударения движущихся частиц эфира, разделённых граничной поверхностью, пространственная форма которой должна соответствовать критерию бесконечной повторяемости физического сценария этих соударений.
\end{abstract}
%\newpage
\vspace{2em}
